\chapter{The Template Library}\label{ch:library}

\file{library/} holds 174 templates plus the full file template, 175 snippets in all, covering the USACO Bronze
to Platinum syllabus and the advanced topics beyond it. The second edition of the handbook added 45 of them (discrete
logarithms, polynomial tools, FWHT, Du's sieve and prime counting, Hungarian and blossom matching, Gomory--Hu and
Stoer--Wagner cuts, the eertree, Lyndon factorization, CDQ, Mo on trees and with updates, persistent treaps, Euler tour
trees, half-planes, rotating calipers and more). Every one of them:

\begin{itemize}
  \item compiles on its own as strict C++17 with zero warnings (\texttt{-Wall -Wextra -Wshadow -Wpedantic}), both as
        written and under \cc{#define int long long};
  \item was checked against a brute-force implementation on hundreds to thousands of random inputs (and planted bugs
        were confirmed to be caught);
  \item avoids every name the template defines (\cc{INF}, \cc{MOD}, \cc{sz}, \cc{all}, \cc{rep}, \cc{dbg},
        \ldots) and uses no \cc{__int128}, so it also works on 32-bit judges.
\end{itemize}

\section{Three ways to insert a template}

\begin{enumerate}
  \item \textbf{Chords.} \key{Ctrl+Alt+T} (core), \key{Ctrl+Alt+Y} (advanced) or \key{Ctrl+Alt+U} (specialist
        techniques), release, then one key. The code is pasted at the cursor and re-indented to match it. The maps below
        show every chord.
  \item \textbf{Prefixes.} Every template has a prefix starting with \texttt{cp}: type it and press \key{Tab}
        (\texttt{cpdsu}, \texttt{cpseg}, \texttt{cphld}). Typing just \texttt{cp} lists all of them with their descriptions.
  \item \textbf{Picker.} \key{Ctrl+Alt+T} then \key{Space} opens VS Code's snippet list; type any part of a name.
\end{enumerate}

\section{Where to paste}

Each template is one of two kinds, shown on its card in Chapter~\ref{ch:reference} and in \file{CHEATSHEET.md}:

\begin{itemize}
  \item \textbf{above solve()} (\emph{global}): a \cc{struct} or functions, such as \cc{DSU}, \cc{SegTree},
        \cc{dijkstra}. Put the cursor on an empty line above \cc{void solve()} and paste; then use it inside
        \cc{solve()}.
  \item \textbf{inside solve()}: a block of statements that uses variables you already have, such as the DFS or grid
        BFS. The first line of the card says which variables it needs (\texttt{needs n, adj}). Paste it where those
        variables exist, then fill in the marked lines.
\end{itemize}

A few templates build on another one and say so (\texttt{Needs cppoint}, \texttt{Needs cpdinic}, \texttt{Needs cpdsurb}):
paste the dependency first.

\section{Conventions used everywhere}

\begin{center}
\begin{tabularx}{\linewidth}{p{4cm}L}
\toprule
\thd{Convention} & \thd{Meaning} \\
\midrule
0-indexed & vertices $0..n-1$, array positions $0..n-1$. \texttt{cpgraph} converts 1-indexed input for you. \\
half-open ranges & every range query is on $[l, r)$: \cc{st.query(0, n)} is the whole array, \cc{query(l, l)} is empty. \\
unreachable & shortest-path templates return \cc{LLONG_MAX} (or the named constant they define) for unreachable vertices. \\
empty result = impossible & e.g.\ \cc{topo_sort} returns an empty vector on a cycle, \cc{TwoSat::solve} when unsatisfiable. \\
\texttt{MOD} & modular templates use the template's \cc{MOD}. \\
edit points & lines to adapt are marked with a comment (\texttt{// transition}, \texttt{// replace with the real check}). \\
\bottomrule
\end{tabularx}
\end{center}

\section{Keyboard maps}

Orange keys have a template; the short name under the key is its prefix without \texttt{cp}. The complete list with
descriptions is Chapter~\ref{ch:reference}.

% generated by tools/cp.py setup -- do not edit
\begin{figure}[H]\centering
\begin{tikzpicture}[kb/.style={draw=cpgray!45, fill=white, rounded corners=2pt, minimum width=11.6mm, minimum height=10.4mm, align=center, font=\sffamily\fontsize{5.6}{6.4}\selectfont, inner sep=0.6pt, text=cpgray!70}, kbused/.style={fill=cporange!10, draw=cporange!70, text=black}]
\node[kb, kbused] at (0.00,0.00) {\textbf{\color{cpblue}\textasciigrave{}}\\xortrie};
\node[kb, kbused] at (1.24,0.00) {\textbf{\color{cpblue}1}\\prefix};
\node[kb, kbused] at (2.48,0.00) {\textbf{\color{cpblue}2}\\bipartite};
\node[kb, kbused] at (3.72,0.00) {\textbf{\color{cpblue}3}\\diff};
\node[kb, kbused] at (4.96,0.00) {\textbf{\color{cpblue}4}\\diff2d};
\node[kb, kbused] at (6.20,0.00) {\textbf{\color{cpblue}5}\\twoptr};
\node[kb, kbused] at (7.44,0.00) {\textbf{\color{cpblue}6}\\window};
\node[kb, kbused] at (8.68,0.00) {\textbf{\color{cpblue}7}\\monostack};
\node[kb, kbused] at (9.92,0.00) {\textbf{\color{cpblue}8}\\subsets};
\node[kb, kbused] at (11.16,0.00) {\textbf{\color{cpblue}9}\\perm};
\node[kb, kbused] at (12.40,0.00) {\textbf{\color{cpblue}0}\\01bfs};
\node[kb, kbused] at (13.64,0.00) {\textbf{\color{cpblue}-}\\rangedp};
\node[kb, kbused] at (14.88,0.00) {\textbf{\color{cpblue}=}\\digit};
\node[kb, kbused] at (0.62,-1.12) {\textbf{\color{cpblue}Q}\\ternary};
\node[kb, kbused] at (1.86,-1.12) {\textbf{\color{cpblue}W}\\floyd};
\node[kb, kbused] at (3.10,-1.12) {\textbf{\color{cpblue}E}\\tree};
\node[kb, kbused] at (4.34,-1.12) {\textbf{\color{cpblue}R}\\sparse};
\node[kb, kbused] at (5.58,-1.12) {\textbf{\color{cpblue}T}\\template};
\node[kb, kbused] at (6.82,-1.12) {\textbf{\color{cpblue}Y}\\bits};
\node[kb, kbused] at (8.06,-1.12) {\textbf{\color{cpblue}U}\\dsu};
\node[kb, kbused] at (9.30,-1.12) {\textbf{\color{cpblue}I}\\bit};
\node[kb, kbused] at (10.54,-1.12) {\textbf{\color{cpblue}O}\\topo};
\node[kb, kbused] at (11.78,-1.12) {\textbf{\color{cpblue}P}\\sieve};
\node[kb, kbused] at (13.02,-1.12) {\textbf{\color{cpblue}[}\\lcs};
\node[kb, kbused] at (14.26,-1.12) {\textbf{\color{cpblue}]}\\tsp};
\node[kb, kbused] at (0.93,-2.24) {\textbf{\color{cpblue}A}\\bsearch};
\node[kb, kbused] at (2.17,-2.24) {\textbf{\color{cpblue}S}\\seg};
\node[kb, kbused] at (3.41,-2.24) {\textbf{\color{cpblue}D}\\dfs};
\node[kb, kbused] at (4.65,-2.24) {\textbf{\color{cpblue}F}\\flood};
\node[kb, kbused] at (5.89,-2.24) {\textbf{\color{cpblue}G}\\gridbfs};
\node[kb, kbused] at (7.13,-2.24) {\textbf{\color{cpblue}H}\\strhash};
\node[kb, kbused] at (8.37,-2.24) {\textbf{\color{cpblue}J}\\dijkstra};
\node[kb, kbused] at (9.61,-2.24) {\textbf{\color{cpblue}K}\\kruskal};
\node[kb, kbused] at (10.85,-2.24) {\textbf{\color{cpblue}L}\\lca};
\node[kb, kbused] at (12.09,-2.24) {\textbf{\color{cpblue};}\\treedp};
\node[kb, kbused] at (13.33,-2.24) {\textbf{\color{cpblue}\textquotesingle{}}\\reroot};
\node[kb, kbused] at (1.55,-3.36) {\textbf{\color{cpblue}Z}\\z};
\node[kb, kbused] at (2.79,-3.36) {\textbf{\color{cpblue}X}\\extgcd};
\node[kb, kbused] at (4.03,-3.36) {\textbf{\color{cpblue}C}\\scc};
\node[kb, kbused] at (5.27,-3.36) {\textbf{\color{cpblue}V}\\read};
\node[kb, kbused] at (6.51,-3.36) {\textbf{\color{cpblue}B}\\bfs};
\node[kb, kbused] at (7.75,-3.36) {\textbf{\color{cpblue}N}\\ncr};
\node[kb, kbused] at (8.99,-3.36) {\textbf{\color{cpblue}M}\\modpow};
\node[kb, kbused] at (10.23,-3.36) {\textbf{\color{cpblue},}\\funcgraph};
\node[kb, kbused] at (11.47,-3.36) {\textbf{\color{cpblue}.}\\jump};
\node[kb, kbused] at (12.71,-3.36) {\textbf{\color{cpblue}/}\\cycle};
\end{tikzpicture}
\caption{Ctrl+Alt+T, then a key no Shift (46 templates)}
\end{figure}
\begin{figure}[H]\centering
\begin{tikzpicture}[kb/.style={draw=cpgray!45, fill=white, rounded corners=2pt, minimum width=11.6mm, minimum height=10.4mm, align=center, font=\sffamily\fontsize{5.6}{6.4}\selectfont, inner sep=0.6pt, text=cpgray!70}, kbused/.style={fill=cporange!10, draw=cporange!70, text=black}]
\node[kb] at (0.00,0.00) {\textbf{\color{cpblue}\textasciigrave{}}};
\node[kb, kbused] at (1.24,0.00) {\textbf{\color{cpblue}1}\\bitsetsum};
\node[kb, kbused] at (2.48,0.00) {\textbf{\color{cpblue}2}\\divsieve};
\node[kb, kbused] at (3.72,0.00) {\textbf{\color{cpblue}3}\\point};
\node[kb, kbused] at (4.96,0.00) {\textbf{\color{cpblue}4}\\rect};
\node[kb, kbused] at (6.20,0.00) {\textbf{\color{cpblue}5}\\prim};
\node[kb] at (7.44,0.00) {\textbf{\color{cpblue}6}};
\node[kb] at (8.68,0.00) {\textbf{\color{cpblue}7}};
\node[kb] at (9.92,0.00) {\textbf{\color{cpblue}8}};
\node[kb] at (11.16,0.00) {\textbf{\color{cpblue}9}};
\node[kb] at (12.40,0.00) {\textbf{\color{cpblue}0}};
\node[kb] at (13.64,0.00) {\textbf{\color{cpblue}-}};
\node[kb] at (14.88,0.00) {\textbf{\color{cpblue}=}};
\node[kb, kbused] at (0.62,-1.12) {\textbf{\color{cpblue}Q}\\knap};
\node[kb, kbused] at (1.86,-1.12) {\textbf{\color{cpblue}W}\\bellman};
\node[kb, kbused] at (3.10,-1.12) {\textbf{\color{cpblue}E}\\graphw};
\node[kb, kbused] at (4.34,-1.12) {\textbf{\color{cpblue}R}\\rng};
\node[kb, kbused] at (5.58,-1.12) {\textbf{\color{cpblue}T}\\trie};
\node[kb, kbused] at (6.82,-1.12) {\textbf{\color{cpblue}Y}\\backtrack};
\node[kb, kbused] at (8.06,-1.12) {\textbf{\color{cpblue}U}\\mitm};
\node[kb, kbused] at (9.30,-1.12) {\textbf{\color{cpblue}I}\\interacti.};
\node[kb, kbused] at (10.54,-1.12) {\textbf{\color{cpblue}O}\\oset};
\node[kb, kbused] at (11.78,-1.12) {\textbf{\color{cpblue}P}\\factor};
\node[kb] at (13.02,-1.12) {\textbf{\color{cpblue}[}};
\node[kb] at (14.26,-1.12) {\textbf{\color{cpblue}]}};
\node[kb, kbused] at (0.93,-2.24) {\textbf{\color{cpblue}A}\\prefix2d};
\node[kb, kbused] at (2.17,-2.24) {\textbf{\color{cpblue}S}\\lazy};
\node[kb, kbused] at (3.41,-2.24) {\textbf{\color{cpblue}D}\\diameter};
\node[kb, kbused] at (4.65,-2.24) {\textbf{\color{cpblue}F}\\bsreal};
\node[kb, kbused] at (5.89,-2.24) {\textbf{\color{cpblue}G}\\graph};
\node[kb, kbused] at (7.13,-2.24) {\textbf{\color{cpblue}H}\\hash};
\node[kb, kbused] at (8.37,-2.24) {\textbf{\color{cpblue}J}\\argsort};
\node[kb, kbused] at (9.61,-2.24) {\textbf{\color{cpblue}K}\\kmp};
\node[kb, kbused] at (10.85,-2.24) {\textbf{\color{cpblue}L}\\lis};
\node[kb] at (12.09,-2.24) {\textbf{\color{cpblue};}};
\node[kb] at (13.33,-2.24) {\textbf{\color{cpblue}\textquotesingle{}}};
\node[kb, kbused] at (1.55,-3.36) {\textbf{\color{cpblue}Z}\\gridpaths};
\node[kb, kbused] at (2.79,-3.36) {\textbf{\color{cpblue}X}\\matrix};
\node[kb, kbused] at (4.03,-3.36) {\textbf{\color{cpblue}C}\\compress};
\node[kb, kbused] at (5.27,-3.36) {\textbf{\color{cpblue}V}\\hull};
\node[kb, kbused] at (6.51,-3.36) {\textbf{\color{cpblue}B}\\bridges};
\node[kb, kbused] at (7.75,-3.36) {\textbf{\color{cpblue}N}\\intervals};
\node[kb, kbused] at (8.99,-3.36) {\textbf{\color{cpblue}M}\\mint};
\node[kb] at (10.23,-3.36) {\textbf{\color{cpblue},}};
\node[kb] at (11.47,-3.36) {\textbf{\color{cpblue}.}};
\node[kb] at (12.71,-3.36) {\textbf{\color{cpblue}/}};
\end{tikzpicture}
\caption{Ctrl+Alt+T, then a key with Shift (31 templates)}
\end{figure}
\begin{figure}[H]\centering
\begin{tikzpicture}[kb/.style={draw=cpgray!45, fill=white, rounded corners=2pt, minimum width=11.6mm, minimum height=10.4mm, align=center, font=\sffamily\fontsize{5.6}{6.4}\selectfont, inner sep=0.6pt, text=cpgray!70}, kbused/.style={fill=cporange!10, draw=cporange!70, text=black}]
\node[kb] at (0.00,0.00) {\textbf{\color{cpblue}\textasciigrave{}}};
\node[kb, kbused] at (1.24,0.00) {\textbf{\color{cpblue}1}\\crt};
\node[kb, kbused] at (2.48,0.00) {\textbf{\color{cpblue}2}\\2sat};
\node[kb, kbused] at (3.72,0.00) {\textbf{\color{cpblue}3}\\incexc};
\node[kb, kbused] at (4.96,0.00) {\textbf{\color{cpblue}4}\\sa};
\node[kb, kbused] at (6.20,0.00) {\textbf{\color{cpblue}5}\\sam};
\node[kb, kbused] at (7.44,0.00) {\textbf{\color{cpblue}6}\\manacher};
\node[kb, kbused] at (8.68,0.00) {\textbf{\color{cpblue}7}\\seginter};
\node[kb, kbused] at (9.92,0.00) {\textbf{\color{cpblue}8}\\polygon};
\node[kb, kbused] at (11.16,0.00) {\textbf{\color{cpblue}9}\\closest};
\node[kb] at (12.40,0.00) {\textbf{\color{cpblue}0}};
\node[kb] at (13.64,0.00) {\textbf{\color{cpblue}-}};
\node[kb] at (14.88,0.00) {\textbf{\color{cpblue}=}};
\node[kb, kbused] at (0.62,-1.12) {\textbf{\color{cpblue}Q}\\dnc};
\node[kb, kbused] at (1.86,-1.12) {\textbf{\color{cpblue}W}\\wavelet};
\node[kb, kbused] at (3.10,-1.12) {\textbf{\color{cpblue}E}\\euler};
\node[kb, kbused] at (4.34,-1.12) {\textbf{\color{cpblue}R}\\dsurb};
\node[kb, kbused] at (5.58,-1.12) {\textbf{\color{cpblue}T}\\treap};
\node[kb, kbused] at (6.82,-1.12) {\textbf{\color{cpblue}Y}\\dinic};
\node[kb, kbused] at (8.06,-1.12) {\textbf{\color{cpblue}U}\\sos};
\node[kb, kbused] at (9.30,-1.12) {\textbf{\color{cpblue}I}\\bitrange};
\node[kb, kbused] at (10.54,-1.12) {\textbf{\color{cpblue}O}\\mo};
\node[kb, kbused] at (11.78,-1.12) {\textbf{\color{cpblue}P}\\persist};
\node[kb] at (13.02,-1.12) {\textbf{\color{cpblue}[}};
\node[kb] at (14.26,-1.12) {\textbf{\color{cpblue}]}};
\node[kb, kbused] at (0.93,-2.24) {\textbf{\color{cpblue}A}\\lazyassign};
\node[kb, kbused] at (2.17,-2.24) {\textbf{\color{cpblue}S}\\beats};
\node[kb, kbused] at (3.41,-2.24) {\textbf{\color{cpblue}D}\\dynseg};
\node[kb, kbused] at (4.65,-2.24) {\textbf{\color{cpblue}F}\\fft};
\node[kb, kbused] at (5.89,-2.24) {\textbf{\color{cpblue}G}\\grundy};
\node[kb, kbused] at (7.13,-2.24) {\textbf{\color{cpblue}H}\\hld};
\node[kb] at (8.37,-2.24) {\textbf{\color{cpblue}J}};
\node[kb, kbused] at (9.61,-2.24) {\textbf{\color{cpblue}K}\\krt};
\node[kb, kbused] at (10.85,-2.24) {\textbf{\color{cpblue}L}\\lineconta.};
\node[kb] at (12.09,-2.24) {\textbf{\color{cpblue};}};
\node[kb] at (13.33,-2.24) {\textbf{\color{cpblue}\textquotesingle{}}};
\node[kb, kbused] at (1.55,-3.36) {\textbf{\color{cpblue}Z}\\monocht};
\node[kb, kbused] at (2.79,-3.36) {\textbf{\color{cpblue}X}\\pollard};
\node[kb, kbused] at (4.03,-3.36) {\textbf{\color{cpblue}C}\\centroid};
\node[kb, kbused] at (5.27,-3.36) {\textbf{\color{cpblue}V}\\virtual};
\node[kb, kbused] at (6.51,-3.36) {\textbf{\color{cpblue}B}\\2ecc};
\node[kb, kbused] at (7.75,-3.36) {\textbf{\color{cpblue}N}\\segnode};
\node[kb, kbused] at (8.99,-3.36) {\textbf{\color{cpblue}M}\\small2lar.};
\node[kb] at (10.23,-3.36) {\textbf{\color{cpblue},}};
\node[kb] at (11.47,-3.36) {\textbf{\color{cpblue}.}};
\node[kb] at (12.71,-3.36) {\textbf{\color{cpblue}/}};
\end{tikzpicture}
\caption{Ctrl+Alt+Y, then a key no Shift (34 templates)}
\end{figure}
\begin{figure}[H]\centering
\begin{tikzpicture}[kb/.style={draw=cpgray!45, fill=white, rounded corners=2pt, minimum width=11.6mm, minimum height=10.4mm, align=center, font=\sffamily\fontsize{5.6}{6.4}\selectfont, inner sep=0.6pt, text=cpgray!70}, kbused/.style={fill=cporange!10, draw=cporange!70, text=black}]
\node[kb] at (0.00,0.00) {\textbf{\color{cpblue}\textasciigrave{}}};
\node[kb] at (1.24,0.00) {\textbf{\color{cpblue}1}};
\node[kb] at (2.48,0.00) {\textbf{\color{cpblue}2}};
\node[kb] at (3.72,0.00) {\textbf{\color{cpblue}3}};
\node[kb] at (4.96,0.00) {\textbf{\color{cpblue}4}};
\node[kb] at (6.20,0.00) {\textbf{\color{cpblue}5}};
\node[kb] at (7.44,0.00) {\textbf{\color{cpblue}6}};
\node[kb] at (8.68,0.00) {\textbf{\color{cpblue}7}};
\node[kb] at (9.92,0.00) {\textbf{\color{cpblue}8}};
\node[kb] at (11.16,0.00) {\textbf{\color{cpblue}9}};
\node[kb] at (12.40,0.00) {\textbf{\color{cpblue}0}};
\node[kb] at (13.64,0.00) {\textbf{\color{cpblue}-}};
\node[kb] at (14.88,0.00) {\textbf{\color{cpblue}=}};
\node[kb, kbused] at (0.62,-1.12) {\textbf{\color{cpblue}Q}\\knuth};
\node[kb] at (1.86,-1.12) {\textbf{\color{cpblue}W}};
\node[kb, kbused] at (3.10,-1.12) {\textbf{\color{cpblue}E}\\phi};
\node[kb, kbused] at (4.34,-1.12) {\textbf{\color{cpblue}R}\\offlinecc};
\node[kb, kbused] at (5.58,-1.12) {\textbf{\color{cpblue}T}\\ntt};
\node[kb, kbused] at (6.82,-1.12) {\textbf{\color{cpblue}Y}\\mcmf};
\node[kb] at (8.06,-1.12) {\textbf{\color{cpblue}U}};
\node[kb, kbused] at (9.30,-1.12) {\textbf{\color{cpblue}I}\\bit2d};
\node[kb] at (10.54,-1.12) {\textbf{\color{cpblue}O}};
\node[kb, kbused] at (11.78,-1.12) {\textbf{\color{cpblue}P}\\profile};
\node[kb] at (13.02,-1.12) {\textbf{\color{cpblue}[}};
\node[kb] at (14.26,-1.12) {\textbf{\color{cpblue}]}};
\node[kb, kbused] at (0.93,-2.24) {\textbf{\color{cpblue}A}\\aho};
\node[kb, kbused] at (2.17,-2.24) {\textbf{\color{cpblue}S}\\slope};
\node[kb, kbused] at (3.41,-2.24) {\textbf{\color{cpblue}D}\\distinct};
\node[kb, kbused] at (4.65,-2.24) {\textbf{\color{cpblue}F}\\floorbloc.};
\node[kb, kbused] at (5.89,-2.24) {\textbf{\color{cpblue}G}\\flowlb};
\node[kb, kbused] at (7.13,-2.24) {\textbf{\color{cpblue}H}\\hopcroft};
\node[kb] at (8.37,-2.24) {\textbf{\color{cpblue}J}};
\node[kb, kbused] at (9.61,-2.24) {\textbf{\color{cpblue}K}\\treeknap};
\node[kb, kbused] at (10.85,-2.24) {\textbf{\color{cpblue}L}\\lichao};
\node[kb] at (12.09,-2.24) {\textbf{\color{cpblue};}};
\node[kb] at (13.33,-2.24) {\textbf{\color{cpblue}\textquotesingle{}}};
\node[kb] at (1.55,-3.36) {\textbf{\color{cpblue}Z}};
\node[kb, kbused] at (2.79,-3.36) {\textbf{\color{cpblue}X}\\xorbasis};
\node[kb, kbused] at (4.03,-3.36) {\textbf{\color{cpblue}C}\\mincount};
\node[kb] at (5.27,-3.36) {\textbf{\color{cpblue}V}};
\node[kb, kbused] at (6.51,-3.36) {\textbf{\color{cpblue}B}\\bcc};
\node[kb, kbused] at (7.75,-3.36) {\textbf{\color{cpblue}N}\\pascal};
\node[kb, kbused] at (8.99,-3.36) {\textbf{\color{cpblue}M}\\mobius};
\node[kb] at (10.23,-3.36) {\textbf{\color{cpblue},}};
\node[kb] at (11.47,-3.36) {\textbf{\color{cpblue}.}};
\node[kb] at (12.71,-3.36) {\textbf{\color{cpblue}/}};
\end{tikzpicture}
\caption{Ctrl+Alt+Y, then a key with Shift (20 templates)}
\end{figure}
\begin{figure}[H]\centering
\begin{tikzpicture}[kb/.style={draw=cpgray!45, fill=white, rounded corners=2pt, minimum width=11.6mm, minimum height=10.4mm, align=center, font=\sffamily\fontsize{5.6}{6.4}\selectfont, inner sep=0.6pt, text=cpgray!70}, kbused/.style={fill=cporange!10, draw=cporange!70, text=black}]
\node[kb] at (0.00,0.00) {\textbf{\color{cpblue}\textasciigrave{}}};
\node[kb] at (1.24,0.00) {\textbf{\color{cpblue}1}};
\node[kb, kbused] at (2.48,0.00) {\textbf{\color{cpblue}2}\\bibfs};
\node[kb, kbused] at (3.72,0.00) {\textbf{\color{cpblue}3}\\cdq};
\node[kb, kbused] at (4.96,0.00) {\textbf{\color{cpblue}4}\\halfplane};
\node[kb, kbused] at (6.20,0.00) {\textbf{\color{cpblue}5}\\calipers};
\node[kb] at (7.44,0.00) {\textbf{\color{cpblue}6}};
\node[kb] at (8.68,0.00) {\textbf{\color{cpblue}7}};
\node[kb] at (9.92,0.00) {\textbf{\color{cpblue}8}};
\node[kb, kbused] at (11.16,0.00) {\textbf{\color{cpblue}9}\\debruijn};
\node[kb, kbused] at (12.40,0.00) {\textbf{\color{cpblue}0}\\gray};
\node[kb] at (13.64,0.00) {\textbf{\color{cpblue}-}};
\node[kb] at (14.88,0.00) {\textbf{\color{cpblue}=}};
\node[kb, kbused] at (0.62,-1.12) {\textbf{\color{cpblue}Q}\\sqrtmod};
\node[kb, kbused] at (1.86,-1.12) {\textbf{\color{cpblue}W}\\fwht};
\node[kb, kbused] at (3.10,-1.12) {\textbf{\color{cpblue}E}\\ett};
\node[kb, kbused] at (4.34,-1.12) {\textbf{\color{cpblue}R}\\ptreap};
\node[kb, kbused] at (5.58,-1.12) {\textbf{\color{cpblue}T}\\disjoints.};
\node[kb, kbused] at (6.82,-1.12) {\textbf{\color{cpblue}Y}\\lyndon};
\node[kb, kbused] at (8.06,-1.12) {\textbf{\color{cpblue}U}\\prufer};
\node[kb, kbused] at (9.30,-1.12) {\textbf{\color{cpblue}I}\\lagrange};
\node[kb, kbused] at (10.54,-1.12) {\textbf{\color{cpblue}O}\\motree};
\node[kb, kbused] at (11.78,-1.12) {\textbf{\color{cpblue}P}\\primecount};
\node[kb] at (13.02,-1.12) {\textbf{\color{cpblue}[}};
\node[kb] at (14.26,-1.12) {\textbf{\color{cpblue}]}};
\node[kb, kbused] at (0.93,-2.24) {\textbf{\color{cpblue}A}\\gomoryhu};
\node[kb, kbused] at (2.17,-2.24) {\textbf{\color{cpblue}S}\\stirling};
\node[kb, kbused] at (3.41,-2.24) {\textbf{\color{cpblue}D}\\dusieve};
\node[kb, kbused] at (4.65,-2.24) {\textbf{\color{cpblue}F}\\cyclefind};
\node[kb, kbused] at (5.89,-2.24) {\textbf{\color{cpblue}G}\\primroot};
\node[kb, kbused] at (7.13,-2.24) {\textbf{\color{cpblue}H}\\treehash};
\node[kb, kbused] at (8.37,-2.24) {\textbf{\color{cpblue}J}\\jumpptr};
\node[kb, kbused] at (9.61,-2.24) {\textbf{\color{cpblue}K}\\kitamasa};
\node[kb, kbused] at (10.85,-2.24) {\textbf{\color{cpblue}L}\\dlog};
\node[kb] at (12.09,-2.24) {\textbf{\color{cpblue};}};
\node[kb] at (13.33,-2.24) {\textbf{\color{cpblue}\textquotesingle{}}};
\node[kb, kbused] at (1.55,-3.36) {\textbf{\color{cpblue}Z}\\anneal};
\node[kb, kbused] at (2.79,-3.36) {\textbf{\color{cpblue}X}\\xorhash};
\node[kb, kbused] at (4.03,-3.36) {\textbf{\color{cpblue}C}\\cartesian};
\node[kb, kbused] at (5.27,-3.36) {\textbf{\color{cpblue}V}\\konig};
\node[kb, kbused] at (6.51,-3.36) {\textbf{\color{cpblue}B}\\sternbroc.};
\node[kb, kbused] at (7.75,-3.36) {\textbf{\color{cpblue}N}\\gauss};
\node[kb, kbused] at (8.99,-3.36) {\textbf{\color{cpblue}M}\\monominima};
\node[kb] at (10.23,-3.36) {\textbf{\color{cpblue},}};
\node[kb] at (11.47,-3.36) {\textbf{\color{cpblue}.}};
\node[kb] at (12.71,-3.36) {\textbf{\color{cpblue}/}};
\end{tikzpicture}
\caption{Ctrl+Alt+U, then a key no Shift (32 templates)}
\end{figure}
\begin{figure}[H]\centering
\begin{tikzpicture}[kb/.style={draw=cpgray!45, fill=white, rounded corners=2pt, minimum width=11.6mm, minimum height=10.4mm, align=center, font=\sffamily\fontsize{5.6}{6.4}\selectfont, inner sep=0.6pt, text=cpgray!70}, kbused/.style={fill=cporange!10, draw=cporange!70, text=black}]
\node[kb] at (0.00,0.00) {\textbf{\color{cpblue}\textasciigrave{}}};
\node[kb] at (1.24,0.00) {\textbf{\color{cpblue}1}};
\node[kb] at (2.48,0.00) {\textbf{\color{cpblue}2}};
\node[kb] at (3.72,0.00) {\textbf{\color{cpblue}3}};
\node[kb] at (4.96,0.00) {\textbf{\color{cpblue}4}};
\node[kb] at (6.20,0.00) {\textbf{\color{cpblue}5}};
\node[kb] at (7.44,0.00) {\textbf{\color{cpblue}6}};
\node[kb] at (8.68,0.00) {\textbf{\color{cpblue}7}};
\node[kb] at (9.92,0.00) {\textbf{\color{cpblue}8}};
\node[kb] at (11.16,0.00) {\textbf{\color{cpblue}9}};
\node[kb] at (12.40,0.00) {\textbf{\color{cpblue}0}};
\node[kb] at (13.64,0.00) {\textbf{\color{cpblue}-}};
\node[kb] at (14.88,0.00) {\textbf{\color{cpblue}=}};
\node[kb] at (0.62,-1.12) {\textbf{\color{cpblue}Q}};
\node[kb] at (1.86,-1.12) {\textbf{\color{cpblue}W}};
\node[kb, kbused] at (3.10,-1.12) {\textbf{\color{cpblue}E}\\eertree};
\node[kb, kbused] at (4.34,-1.12) {\textbf{\color{cpblue}R}\\radix};
\node[kb, kbused] at (5.58,-1.12) {\textbf{\color{cpblue}T}\\sack};
\node[kb] at (6.82,-1.12) {\textbf{\color{cpblue}Y}};
\node[kb] at (8.06,-1.12) {\textbf{\color{cpblue}U}};
\node[kb, kbused] at (9.30,-1.12) {\textbf{\color{cpblue}I}\\diffcons};
\node[kb, kbused] at (10.54,-1.12) {\textbf{\color{cpblue}O}\\moupdates};
\node[kb, kbused] at (11.78,-1.12) {\textbf{\color{cpblue}P}\\poly};
\node[kb] at (13.02,-1.12) {\textbf{\color{cpblue}[}};
\node[kb] at (14.26,-1.12) {\textbf{\color{cpblue}]}};
\node[kb, kbused] at (0.93,-2.24) {\textbf{\color{cpblue}A}\\stoerwagn.};
\node[kb, kbused] at (2.17,-2.24) {\textbf{\color{cpblue}S}\\subsetconv};
\node[kb, kbused] at (3.41,-2.24) {\textbf{\color{cpblue}D}\\det};
\node[kb] at (4.65,-2.24) {\textbf{\color{cpblue}F}};
\node[kb, kbused] at (5.89,-2.24) {\textbf{\color{cpblue}G}\\blossom};
\node[kb, kbused] at (7.13,-2.24) {\textbf{\color{cpblue}H}\\hungarian};
\node[kb] at (8.37,-2.24) {\textbf{\color{cpblue}J}};
\node[kb] at (9.61,-2.24) {\textbf{\color{cpblue}K}};
\node[kb, kbused] at (10.85,-2.24) {\textbf{\color{cpblue}L}\\offlinelca};
\node[kb] at (12.09,-2.24) {\textbf{\color{cpblue};}};
\node[kb] at (13.33,-2.24) {\textbf{\color{cpblue}\textquotesingle{}}};
\node[kb] at (1.55,-3.36) {\textbf{\color{cpblue}Z}};
\node[kb] at (2.79,-3.36) {\textbf{\color{cpblue}X}};
\node[kb, kbused] at (4.03,-3.36) {\textbf{\color{cpblue}C}\\minplus};
\node[kb] at (5.27,-3.36) {\textbf{\color{cpblue}V}};
\node[kb, kbused] at (6.51,-3.36) {\textbf{\color{cpblue}B}\\bm};
\node[kb, kbused] at (7.75,-3.36) {\textbf{\color{cpblue}N}\\longpath};
\node[kb, kbused] at (8.99,-3.36) {\textbf{\color{cpblue}M}\\mincircle};
\node[kb] at (10.23,-3.36) {\textbf{\color{cpblue},}};
\node[kb] at (11.47,-3.36) {\textbf{\color{cpblue}.}};
\node[kb] at (12.71,-3.36) {\textbf{\color{cpblue}/}};
\end{tikzpicture}
\caption{Ctrl+Alt+U, then a key with Shift (16 templates)}
\end{figure}


\section{Categories}

\begin{center}
\begin{tabularx}{\linewidth}{p{3.4cm}L}
\toprule
\thd{Folder} & \thd{Contents} \\
\midrule
\file{basics/} & input readers, prefix sums, difference arrays, two pointers, sliding window, monotonic stack, binary and
  ternary search, subsets, permutations, backtracking, argsort, intervals, bit helpers, meet in the middle,
  coordinate compression \\
\file{graph/} & DFS, BFS, flood fill, grid BFS, Dijkstra, 0-1 BFS, Bellman--Ford, Floyd--Warshall, topological sort,
  DSU (+ rollback), Kruskal, Prim, SCC, 2-SAT, bipartite check, bridges and cut vertices, 2-edge-connected components,
  biconnected components, Euler paths, cycle finding, functional graphs, k-th successor, offline connectivity \\
\file{tree/} & tree DFS / Euler tour, diameter, LCA, tree DP, rerooting, small-to-large, HLD, centroid decomposition,
  virtual tree, Kruskal reconstruction tree, tree knapsack \\
\file{ds/} & segment trees (point, lazy add, lazy assign, custom node with binary search, min + count, dynamic,
  persistent, beats), Fenwick trees (point, range, 2D), sparse table, ordered set, anti-hack hash, wavelet tree,
  treap, Mo's algorithm, offline distinct count \\
\file{dp/} & knapsacks, bitset subset sum, grid paths, LIS, LCS / edit distance, bitmask (TSP), range DP, digit DP,
  SOS DP, divide-and-conquer and Knuth optimisations, broken profile, LineContainer, Li Chao tree, slope trick \\
\file{math/} & modular power, Mint, nCr, sieve, factorization, Pollard rho, divisor sieve, totient, M\"obius,
  floor blocks, Pascal, extended gcd, CRT, matrix power, FFT, NTT, XOR basis, Grundy numbers, inclusion--exclusion \\
\file{string/} & hashing, KMP, Z-function, trie, binary trie, Aho--Corasick, suffix array + LCP, suffix automaton, Manacher \\
\file{geometry/} & point primitives, segment intersection, polygons (area, containment, Pick), convex hull, closest
  pair, rectangles \\
\file{flow/} & Dinic + min cut, min-cost flow, flows with lower bounds, Hopcroft--Karp \\
\file{misc/} & random numbers, interactive-problem helpers \\
\bottomrule
\end{tabularx}
\end{center}
